
\medskip
\noindent \SGARefTarget{sga7.sga7i.eviii.n5-7}{5.7}.  Let $S$ be a scheme, $T$ a closed subscheme, and $Z$
a set.  The \underline{constant scheme over $S$, truncated along $T$,
with value $Z$}, denoted $Z_{S,T}$, is the scheme obtained by taking
the sum of copies $S_z$ of $S$, one for each $z\in Z$, and gluing all
these copies along the open subscheme $U=S\setminus T$.  Thus
\begin{equation}\tag{\SGARefTarget{sga7.sga7i.eviii.n5-7-1}{5.7.1}}
 Z_{S,T}\times_S U\simeq U,
 \qquad
 Z_{S,T}\times_S T\simeq Z_T,
\end{equation}
and $Z_{S,T}$ is etale over $S$.  For a variable $S$-scheme $S'$
there is an isomorphism, functorial in $S'$,
\begin{equation}\tag{\SGARefTarget{sga7.sga7i.eviii.n5-7-2}{5.7.2}}
 \Hom_S(S',Z_{S,T})\simeq\Hom_T(S'_T,Z_T).
\end{equation}
The map is restriction to $S'_T$, using the second isomorphism in
(\SGARefLink{sga7.sga7i.eviii.n5-7-1}{5.7.1}); we leave it to the reader to check that the map is bijective.
If
\[
 i:T\longrightarrow S
\]
denotes the inclusion, the functorial isomorphism (\SGARefLink{sga7.sga7i.eviii.n5-7-2}{5.7.2}) may also be
written as an isomorphism of sheaves
\begin{equation}\tag{\SGARefTarget{sga7.sga7i.eviii.n5-7-3}{5.7.3}}
 Z_{S,T}\simeq i_*(Z_T).
\end{equation}
This formula, or the definition directly, shows that as $Z$ varies,
$Z_{S,T}$ depends functorially on $Z$, and that this functor commutes
with finite inverse limits.  In particular, it takes groups to group
schemes and commutative groups to commutative group schemes, as is also
clear from the form (\SGARefLink{sga7.sga7i.eviii.n5-7-2}{5.7.2}) of the functor represented by $Z_{S,T}$.

It is also immediate from the definition, or from (\SGARefLink{sga7.sga7i.eviii.n5-7-3}{5.7.3}), that the
formation of $Z_{S,T}$ commutes with every base change
$f:S'\to S$.  This gives an isomorphism, functorial in $Z$,
\begin{equation}\tag{\SGARefTarget{sga7.sga7i.eviii.n5-7-4}{5.7.4}}
 Z_{S,T}\times_S S'\simeq Z_{S',T'},
 \qquad T'=T\times_S S'=T_{S'}.
\end{equation}

\medskip
\noindent \SGARefTarget{sga7.sga7i.eviii.n5-8}{5.8}.  We shall study torsors $E$ under a group scheme of the
form $Z_{S,T}$, where $Z$ is a group, assumed commutative here for
simplicity.  There is a natural functor
\begin{equation}\tag{\SGARefTarget{sga7.sga7i.eviii.n5-8-0}{5.8.0}}
 E\longmapsto E_T
\end{equation}
from the category of torsors under $Z_{S,T}$ to the category of
torsors under $Z_T$.  We claim that this functor is an
\underline{equivalence of categories}.  Its full faithfulness means
that for two torsors $E,F$ under $Z_{S,T}$, the map
\[
 \Hom(E,F)\longrightarrow\Hom(E_T,F_T)
\]
is bijective.  By fpqc descent this reduces immediately to the case in
which $E$ and $F$ are trivial, hence to $E=F=Z_{S,T}$, where the
assertion is precisely the bijectivity of (\SGARefLink{sga7.sga7i.eviii.n5-7-2}{5.7.2}) for $S'=S$.

It remains to prove that (\SGARefLink{sga7.sga7i.eviii.n5-8-0}{5.8.0}) is essentially surjective.  Let $E_T$
be a torsor under $Z_T$.  We already observed in the proof of \SGARefLink{sga7.sga7i.eviii.n5-1}{5.1} that
$E_T$ is representable and locally trivial for the etale topology.  It
therefore suffices to prove that
\[
 H^1(S_{\mathrm{\acute{e}t}},Z_{S,T})
 \longrightarrow H^1(T_{\mathrm{\acute{e}t}},Z_T)
\]
is bijective.  This follows from (\SGARefLink{sga7.sga7i.eviii.n5-7-3}{5.7.3}) and the relation
\[
 R^1i_{\mathrm{\acute{e}t},*}(Z_T)=0,
\]
which is the special case SGA 4 VIII 5.6.

\medskip
\noindent \SGARefTarget{sga7.sga7i.eviii.n5-8-1}{5.8.1}.  Let now $P$ be a commutative group scheme over $S$.
We wish to determine the category of extensions, as sheaves of
commutative groups, of $P$ by $Z_{S,T}$.  For this we use the
description in VII \SGARefLink{sga7.sga7i.evii.n1-1-6}{1.1.6} and \SGARefLink{sga7.sga7i.evii.n1-3}{1.3} in terms of torsors on
$P$, $P\times P$, and $P\times P\times P$, under the inverse images
of $Z_{S,T}$.  By (\SGARefLink{sga7.sga7i.eviii.n5-7-4}{5.7.4}), these inverse images are again of the same
form as $Z_{S,T}$.  Applying the preceding result to those torsors
shows that they correspond to torsors on
$P_T$, $P_T\times_TP_T$, and
$P_T\times_TP_T\times_TP_T$ under the inverse images of $Z_T$.
We therefore obtain the following.

\medskip
\noindent\underline{Proposition \SGARefTarget{sga7.sga7i.eviii.n5-9}{5.9}.}  \underline{Let $S$ be a scheme,
$T$ a closed subscheme, $Z$ a commutative group, $Z_{S,T}$ the
corresponding truncated constant group scheme over $S$ defined in
\SGARefLink{sga7.sga7i.eviii.n5-7}{5.7}, and $P$ a commutative group scheme over $S$.  Then restriction to
$T$ induces an equivalence between the category of sheaves of
commutative groups which are extensions of $P$ by $Z_{S,T}$ and the
category of sheaves of commutative groups which are extensions of
$P_T$ by $Z_T$.}

An essentially identical argument gives the following.

\medskip
\noindent\underline{Corollary \SGARefTarget{sga7.sga7i.eviii.n5-10}{5.10}.}  \underline{With the notation of
\SGARefLink{sga7.sga7i.eviii.n5-9}{5.9}, let $Q$ be another commutative group scheme over $S$.  Then
restriction to $T$ induces an equivalence between the category of
biextensions of $(P,Q)$ by $Z_{S,T}$ and the category of biextensions
of $(P_T,Q_T)$ by $Z_T$.}

